\subsection{A quantitative bound for the original tensor projections}
\label{archimedean-tensor-projections}

The first-length Bessel recursion needs a bound for projections in an
indefinite tensor product.  An arbitrary such projection can have a
large norm.  The parameter separation below controls all the closely
spaced compact parameters together; only the remaining $m$ parameters
enter the exponent.  This proves the original algebraic projection
estimate used by that recursion \cite{tensor-projection-source}.

Let $n\ge2$, $1\le k\le n-1$, $m=n+1-k$, and $B\ge10^6n^7$.
Number coordinates from $0$ to $n-1$, and put
\[
 C=\{0,\ldots,k-2\},\qquad N=\{k-1,\ldots,n-1\}.
\]
Thus $C$ may be empty and $|N|=m$.  Let $s_a\in\{1,-1\}$, with
$s_a=1$ for $a\in C$, and let real parameters $\mu_a$ satisfy
\begin{equation}\label{tensor-projection-shape}
 \begin{split}
 \mu_a-\mu_b&\ge2 &&(a<b),\\
 -B/100\le\mu_a&\le B/8+B/100 &&(a\in C),\\
 -2B-B/100\le\mu_a&\le-B+B/100 &&(a\in N).
 \end{split}
\end{equation}
Let $H$ be a unitary module on a common smooth domain, with matrix
units $E_{ab}$ satisfying the usual commutator relations and
$E_{ab}^*=s_as_bE_{ba}$.  Put
\[
 \rho_a=(n-1)/2-a,\qquad
 (\rho_K)_a=\frac12\bigl(\#\{b>a:s_b=s_a\}
                         -\#\{b<a:s_b=s_a\}\bigr),
 \quad \Lambda=\mu+\rho-2\rho_K.
\]
The full Casimir acts by
$c_G=\|\mu\|^2-\|\rho\|^2$.  Set
$\kappa=\|\Lambda+\rho_K\|^2-\|\rho_K\|^2$.
The weight $\Lambda$ is dominant within each compact block, so
$\kappa\ge\|\Lambda\|^2\ge0$.

A vector has level at most $j$ if it belongs to a finite-dimensional
compact-invariant space on which
\[
              \Omega_K\le(\sqrt\kappa+4nj)^2 I
\]
and is a finite sum of joint weight vectors of weights $\omega$ with
$\|\omega-\Lambda\|_1\le j$.  This is the original level condition.
For fixed $j$ these vectors form a subspace: add the finite compact
witnesses and use the self-adjoint compact Casimir decomposition.

Give $H^n=(\mathbb C^n)^*\otimes H$ its positive Hilbert norm and
the separate indefinite form
\[
             [x,y]=\sum_a s_a\langle x_a,y_a\rangle.
\]
Its tensor action is
$(E^{\rm ten}_{ab}x)_c=E_{ab}x_c-\delta_{bc}x_a$.
Retain the original tensor decomposition: unitary modules $H_t$ of
parameters $\mu-e_t$ and intertwining maps
\[
 i_t:H_t\longrightarrow H^n,\qquad p_t:H^n\longrightarrow H_t,
 \quad p_ti_u=\delta_{tu}I,
 \quad\sum_ti_tp_t=I.
\]
Their images are pairwise orthogonal for $[\ ,\ ]$, and
$[i_ty,i_ty]=\epsilon_t\|y\|^2$ with $\epsilon_t\in\{1,-1\}$.
All maps preserve the smooth domains.  No positive isometry of $i_t$
is assumed.

\begin{theorem}\label{original-tensor-projection-bound}
For a vector $v$ of level at most $j\le4n^2$, and every coordinate $c$,
each $p_t(e_c\otimes v)$ has level at most $j+2$ in $H_t$, and
\[
             \sum_t\|p_t(e_c\otimes v)\|^2
                              \le(6B)^{2m}\|v\|^2.
\]
\end{theorem}

\begin{proof}
We first obtain the small off-diagonal estimate which the original
coarse derivative bound does not by itself supply.  If $u$ has level
$J\le4n^2+2$, the Casimir energy identity gives
\[
 \sum_{a\ne b}\|E_{ab}u\|^2
   =2\langle u,\Omega_Ku\rangle-c_G\|u\|^2
                                      -\sum_a\|E_{aa}u\|^2.
\]
Joint weight spaces are orthogonal, and hence
$\sum_a\|E_{aa}u\|^2
\ge(\|\Lambda\|^2-2\|\Lambda\|_\infty J)\|u\|^2$.
The large quadratic parameter terms cancel exactly:
\begin{equation}\label{tensor-projection-cancellation}
 2\kappa-c_G-\|\Lambda\|^2
              =2\langle\mu,\rho\rangle
                         +2\|\rho\|^2-4\|\rho_K\|^2.
\end{equation}
Here $|\rho_a|,|(\rho_K)_a|\le n/2$, $|\mu_a|\le3B$,
$\|\Lambda\|_\infty\le4B$, $\sqrt\kappa\le4nB$, and
$J\le6n^2$.  Consequently
\[
 \begin{split}
 \sum_{a\ne b}\|E_{ab}u\|^2
 &\le\bigl((384n^4+51n^2)B+1152n^6+n^3/2\bigr)\|u\|^2\\
 &\le512n^4B\|u\|^2.
 \end{split}
\]
Since $B\ge10^6n^7\ge131072n^5$, this proves the original estimate
\begin{equation}\label{tensor-projection-small-offdiagonal}
          n\sum_{a\ne b}\|E_{ab}u\|^2
                                \le(B/16)^2\|u\|^2.
\end{equation}
Thus we reuse the Casimir identity, but require its explicit
cancellation against the diagonal energy.

Write $P_t=i_tp_t$.  The compact Casimir on the standard or dual
standard module is at most $n$.  Tensoring therefore adds at most
$\sqrt n$ to the square root of a compact energy bound.  A compact
intertwiner carries only Casimir eigenvalues already occurring in its
source, whether or not the intertwiner is an isometry.  It follows
that $p_t(e_c\otimes v)$ has compact energy bounded by
$(\sqrt\kappa+4nj+\sqrt n)^2$.

The shifted lowest weight is $\Lambda-e_t$, still compact dominant.
Let its compact scalar be $\kappa_t$.  For $a=\Lambda-e_t$,
$\kappa_t\ge\|a\|^2$ and
$\kappa-\kappa_t=2a_t+1+2(\rho_K)_t
\le2\sqrt{\kappa_t}+n$.
Therefore $\sqrt\kappa\le\sqrt{\kappa_t}+\sqrt n$.
The weights in $p_t(e_c\otimes v)$ are $\omega-e_c$, whose distances
from $\Lambda-e_t$ are at most $j+2$.  This proves its level assertion.
Applying the inclusion and contracting with a standard coordinate
adds a second $\sqrt n$ to the compact bound and changes the weight
to $\omega-e_c+e_a$.  Thus every component of $P_t(e_c\otimes v)$
has level at most $j+2$ in $H$ as well.

Define the split Casimir by
\[
                    (Tx)_a=-\sum_b E_{ab}x_b.
\]
It is self-adjoint for the indefinite form.  Expanding the tensor
Casimir gives $\Omega^{\rm ten}=c_G I+nI+2T$.  On the summand of
parameter $\mu-e_t$ this proves
\[
                   TP_t=\zeta_tP_t,\qquad
                   \zeta_t=-\mu_t+(1-n)/2.
\]
In particular the eigenvalues are distinct and
$|\zeta_t-\zeta_u|\ge2$ for $t\ne u$.

Fix $x=e_c\otimes v$ and let $V$ be the finite span of $P_tx$.
It is $T$-invariant, and all its components have level at most $j+2$.
On such vectors write $T=D+O$, where $(Dy)_a=-E_{aa}y_a$.
Cauchy--Schwarz in each row and
\eqref{tensor-projection-small-offdiagonal} give
\begin{equation}\label{tensor-projection-perturbation}
                     \|Oy\|\le B\|y\|/16\qquad(y\in V).
\end{equation}
The shape inequalities and
$j+2+3n/2\le17B/800$ show that the joint weight values of $D$
belong to
\[
          [-5B/32,B/32]\quad\hbox{on coordinates in }C,
 \qquad [31B/32,65B/32]\quad\hbox{on coordinates in }N.
\]
The same two intervals contain respectively the $\zeta_t$ with
$t\in C,N$.  These are bounds on finite sums of orthogonal weight
vectors; $D$ need not preserve $V$.

On the circle $|z+B/16|=9B/16$, the distance to both intervals is
at least $15B/32$.  Thus
\[
 \|(z-D)y\|\ge15B\|y\|/32,\quad
 \|(z-T)y\|\ge13B\|y\|/32,\quad
                            \|Ty\|\le3B\|y\|.
\]
The resolvent of $T|_V$ exists on that circle.  Integrating its finite
eigenspace expansion gives the sum $P_C=\sum_{t\in C}P_t$ and
\begin{equation}\label{tensor-projection-cloud-bound}
                         \|P_C|_V\|\le18/13.
\end{equation}
This contour identity also follows by integrating the scalar functions
$(z-\zeta_t)^{-1}$, equal to one for the enclosed eigenvalues and zero
for the others.  Hence $\|(I-P_C)|_V\|<3$.
For each $t\in N$, polynomial interpolation now uses only those
$m$ exceptional parameters:
\[
 P_t|_V=
 \prod_{u\in N\setminus\{t\}}
          \frac{T-\zeta_u I}{\zeta_t-\zeta_u}(I-P_C)|_V,
 \qquad \|P_t|_V\|\le3(6B)^{m-1}.
\]

It remains to use the indefinite form without summing all the compact
parameters separately.  If $y=P_tx$ and $t\in C$, then on its $N$
coordinates the eigenvalue equation and the gap of $15B/16$ imply
\[
        \|y_N\|\le\frac{16}{15B}\|Oy\|\le\|y\|/15.
\]
All negative signs lie among these coordinates.  Therefore
$[y,y]\ge(1-2/225)\|y\|^2\ge0$.  Orthogonality of the summands gives
\[
 \begin{split}
 \sum_t\|p_tx\|^2
 &=\sum_t|[P_tx,P_tx]|\\
 &=[P_Cx,P_Cx]+\sum_{t\in N}|[P_tx,P_tx]|\\
 &\le\bigl((18/13)^2+9m(6B)^{2m-2}\bigr)\|x\|^2\\
 &\le(6B)^{2m}\|x\|^2.
 \end{split}
\]
The final inequality uses $1\le m\le n\le B$ and
$9B+2\le36B^2$.  The argument includes $C=\varnothing$ and $v=0$.
Since $\|e_c\otimes v\|=\|v\|$, the theorem follows.
\end{proof}

In the Bessel application $n=2r+k+1$, so $m=2r+2$.  The theorem
therefore supplies exactly the projection factor
$(6B)^{2(2r+2)}$ and the level increase by two in the original
first-length recursion.  Existence of the original tensor
decomposition, its normalized Bessel lowering scalar, and the original
mass/isometry input remain separate obligations; the projection
estimate does not assume or prove those analytic identities.
